\documentclass[review,12pt]{elsarticle}

\usepackage[utf8]{inputenc}
\usepackage[T1]{fontenc}
\usepackage{amsmath,amssymb}
\usepackage{graphicx}
\usepackage{booktabs}
\usepackage{caption}
\usepackage{subcaption}
\usepackage{array}
\usepackage{multirow}
\usepackage[hyphens]{url}
\usepackage{hyperref}
\usepackage{lineno}
\usepackage{setspace}

\hypersetup{colorlinks=true,linkcolor=blue,urlcolor=blue,citecolor=blue}
\doublespacing
\linenumbers

% Gene names in italics
\newcommand{\gene}[1]{\textit{#1}}

\journal{Machine Learning with Applications}

\begin{document}

\begin{frontmatter}
\sloppy

\title{A Biologically Informed Explainable Machine Learning Framework
for Breast Cancer Progression Risk Stratification from Tumor
Gene Expression}

\author[1]{Suhaan Thayyil\corref{cor1}}
\ead{thayyilsuhaan@gmail.com}

\cortext[cor1]{Corresponding author}

\affiliation[1]{organization={Independent Researcher},
               city={Waxhaw},
               state={NC},
               country={USA}}

\begin{abstract}
Gene expression profiling offers a principled supplement to traditional
breast cancer staging, yet most transcriptomic prediction models sacrifice
clinical transparency for discriminative performance.  We present a
biologically informed explainable machine learning framework that
reduces RNA-sequencing data to seven curated pathway activity scores and
evaluates them through both Cox proportional hazards survival analysis
and binary classification, with systematic comparisons against standard
pathway-scoring methods, molecular subtype baselines, and controlled
feature ablation.

Seven pathway scores---covering proliferation, estrogen response, immune
infiltration, apoptosis, epithelial-to-mesenchymal transition, HER2
signalling, and angiogenesis---were derived from the TCGA-BRCA cohort
(\(n = 213\)) and validated externally on SCAN-B GSE96058 (\(n = 1{,}483\)).
Three classifiers (elastic net, random forest, gradient boosting) and a
Cox model were compared under stratified five-fold cross-validation.
SHAP analysis produced per-patient feature attributions, whose stability
was assessed through cross-fold rank correlations.  The mean z-score
aggregation was benchmarked against a rank-based ssGSEA baseline and a
PAM50 molecular subtype-only baseline.

The combined model achieved AUC~0.856 (random forest) and Cox
concordance index~0.827, substantially outperforming the subtype-only
baseline (AUC~0.613; $\Delta$AUC~$+$0.243; $\Delta$C-index~$+$0.214).
Within Luminal~A tumours alone (\(n = 719\)), the model separated
high-risk from low-risk patients with log-rank
\(p = 5.9 \times 10^{-22}\) and cross-validated C-index~0.848, confirming
independent prognostic signal beyond subtype classification.  Mean
z-score aggregation outperformed ssGSEA by 0.010--0.038 AUC across all
classifiers.  SHAP rankings were stable across folds (Spearman
\(\rho = 0.82\)) and recovered established oncological mechanisms.
Decision curve analysis demonstrated positive net benefit across
threshold probabilities 0.04--0.50, and the net reclassification index
versus the subtype baseline was 0.801 (95\% CI: 0.679--0.911).

These results establish that pathway-oriented explainable ML delivers
transparent, generalisable breast cancer risk stratification with
prognostic value beyond molecular subtype classification, and that simple
mean z-score aggregation outperforms the standard ssGSEA approach for
compact curated gene sets.
\end{abstract}

\begin{keyword}
\sloppy
breast cancer prognosis \sep
gene expression \sep
biological pathway scores \sep
explainable AI \sep
SHAP \sep
Cox proportional hazards \sep
molecular subtype baseline \sep
TCGA \sep
SCAN-B
\end{keyword}

\end{frontmatter}

%=======================================================================
\section{Introduction}
\label{sec:intro}
%=======================================================================

Breast cancer is not a single disease but a family of molecularly
distinct subtypes.  At the gene expression level, tumours are commonly
classified into at least four intrinsic categories: Luminal~A,
Luminal~B, HER2-enriched, and Basal-like~\citep{tcga2012}.  This
classification guides treatment decisions because each subtype carries a
different expected trajectory.  Luminal~A tumours tend to express low
Ki-67 and follow a relatively indolent course, while Basal-like (triple-
negative) tumours are characterised by high Ki-67, aggressive growth, and
shorter survival~\citep{dowsett2011}.  However, even within a single
subtype, patient outcomes vary widely, and the subtype label alone leaves
substantial prognostic uncertainty.

Standard clinical variables such as tumour size, nodal involvement,
histological grade, and patient age explain part of this residual
variation.  RNA sequencing offers a much more granular window into
tumour biology, measuring expression across 20,000 or more genes
simultaneously.  This richness, however, creates a challenge: models
that ingest thousands of raw gene-level features are vulnerable to
overfitting, particularly when training cohorts are modest, and they
produce predictions that are difficult for clinicians to interpret.

\subsection{The interpretability gap}
\label{sec:gap}

Commercially available multigene assays such as Oncotype~DX~\citep{paik2004}
and MammaPrint~\citep{vantveer2002,cardoso2016} have demonstrated that
transcriptomic information can meaningfully inform treatment decisions.
These panels are proprietary, fixed in composition, and return a single
risk score without exposing the internal logic.  At the other extreme,
research models built on genome-wide features may achieve high
discrimination but function as opaque black boxes whose predictions
cannot be traced back to recognisable biological processes~\citep{rudin2019}.

A growing body of work has applied explainable artificial intelligence
methods to cancer genomics.  Chakraborty et al.~\citep{chakraborty2021}
used XGBoost with SHAP to relate immune cell composition to five-year
breast cancer survival.  Nkurunziza et al.~\citep{nature2025shap} applied
SHAP to genome-wide features for breast cancer diagnosis.
Chen et al.~\citep{chen2021pathway} constructed a sub-pathway activity matrix
from KEGG annotations for survival stratification.  Standard pathway-
scoring methods such as single-sample gene set enrichment analysis
(ssGSEA)~\citep{barbie2009} and Gene Set Variation Analysis
(GSVA)~\citep{hanzelmann2013} are widely used but produce scores that are
not directly optimised for supervised prediction and have not been
formally benchmarked against simpler aggregation strategies in prognostic
contexts.

Despite this progress, a gap remains.  Most studies either operate at
the individual gene level, sacrificing interpretability, or employ
unsupervised enrichment methods that are not integrated with supervised
modelling.  Few studies benchmark against a molecular subtype-only
baseline, even though well-known clinical markers such as Ki-67,
ER status, and age are strongly associated with PAM50 subtype.  To our
knowledge no published study has combined (1)~manually curated hallmark
pathway scores as the primary supervised feature set, (2)~both survival
analysis and binary classification on the same features, (3)~formal
benchmarking against ssGSEA and molecular subtype baselines,
(4)~per-patient SHAP decomposition, (5)~controlled feature ablation,
(6)~cross-fold feature stability analysis, and (7)~external validation
on a large independent cohort, all within a unified framework.

\subsection{Contributions}
\label{sec:contributions}

This study makes three principal contributions.

\begin{enumerate}
\item \textbf{A biologically informed dimensionality-reduction strategy
      for prognosis.}  Collapsing RNA-seq data into seven pathway scores
      via simple mean z-score aggregation produces a compact feature
      representation that outperforms both the standard ssGSEA baseline
      and a molecular subtype-only classifier (AUC~0.856 vs.~0.613), and
      carries independent prognostic signal in both survival analysis and
      binary classification.

\item \textbf{A generalisable ML pipeline validated across cohorts and modelling paradigms.}
      Cox proportional hazards, elastic net, random forest, and gradient boosting are evaluated
      with systematic ablation and stability analyses
      on two independent cohorts (total $n=1{,}696$).

\item \textbf{Patient-level SHAP interpretability at pathway resolution.}
      Per-patient decompositions ground explanations in recognised hallmarks.
      Within-subtype stratification confirms the framework captures risk
      variation beyond what subtype labels resolve
      (Luminal~A log-rank \(p = 5.9 \times 10^{-22}\)).
\end{enumerate}

\subsection{Scope and non-claims}
\label{sec:scope}

This work is a retrospective proof-of-concept study.  The models
described here are not intended as diagnostic tools and have not been
validated in a prospective clinical setting.  All analyses are
observational, and no causal claims regarding pathway activity and
patient outcomes are made.

%=======================================================================
\section{Related work}
\label{sec:related}
%=======================================================================

Table~\ref{tab:comparison} contextualises the present framework against
recent studies that apply machine learning to breast cancer gene
expression data.

\begin{table}[!ht]
\centering
\caption{Comparison with related breast cancer prediction approaches.
         Columns: Feat.~= feature type; Surv.~= survival analysis;
         XAI = explainability; Abl.~= ablation; Stab.~= stability;
         Sub.~= subtype baseline; Ext.~= external validation.}
\label{tab:comparison}
\resizebox{\linewidth}{!}{%
\begin{tabular}{llllccccc}
\toprule
\textbf{Study} & \textbf{Features} & \textbf{Surv.} &
\textbf{XAI} & \textbf{Abl.} & \textbf{Stab.} &
\textbf{Sub.} & \textbf{Ext.} \\
\midrule
Chakraborty et al.~\citep{chakraborty2021} & Immune & No & SHAP & No & No & No & No \\
Chen et al.~\citep{chen2021pathway}        & KEGG   & Cox & No  & No & No & No & GEO \\
Nkurunziza et al.~\citep{nature2025shap}   & 3,602 genes & No & SHAP & No & No & No & No \\
\textbf{Ours}                              & \textbf{7 pathways} & \textbf{Cox} & \textbf{SHAP} & \textbf{Yes} & \textbf{Yes} & \textbf{Yes} & \textbf{SCAN-B} \\
\bottomrule
\end{tabular}}
\end{table}

Chakraborty et al.\ used XGBoost with SHAP on immune cell composition
estimates from TCGA but did not include survival modelling, feature
ablation, stability analysis, or external validation.
Chen et al.\ applied Cox regression to KEGG pathway features and
validated on GEO, but provided no explainability, ablation, or subtype
benchmarking.  Nkurunziza et al.\ trained classifiers on 3,602 individual
genes with SHAP explanations but lacked survival endpoints and external
cohorts.  The present framework is the only study in this comparison to
address all seven design criteria simultaneously.

%=======================================================================
\section{Materials and methods}
\label{sec:methods}
%=======================================================================

\subsection{Ethical statement and data sources}
\label{sec:data}

All data are publicly available and fully de-identified.  No new
specimens were collected.  This study qualifies as exempt from
Institutional Review Board review under 45~CFR~46.102(e).

\textbf{Training cohort (TCGA-BRCA).}
213 patients with RNA-seq expression (log$_2$ RSEM-normalised) and
clinical annotation from The Cancer Genome Atlas~\citep{tcga2012},
accessed via the GDC Data Portal.

\textbf{Validation cohort (GSE96058 / SCAN-B).}
1,483 patients with RNA-seq expression and overall survival data from
the Sweden Cancerome Analysis Network~\citep{saal2015,brueffer2018},
accessed via NCBI GEO (accession GSE96058).  No GSE96058 data entered
model training or hyperparameter selection.

Table~\ref{tab:cohorts} summarises the two cohorts.

\begin{table}[!ht]
\centering
\caption{Cohort characteristics.  Age and tumour size are reported as
  median (interquartile range).  Percentages are of patients with
  non-missing values.  Events~$=$~deaths; survival endpoint is
  overall survival.  Demographic estimates are derived from published
  summary statistics~\citep{tcga2012,saal2015,brueffer2018}; exact
  per-cohort values should be verified via
  \texttt{notebooks/01\_data\_preprocessing.ipynb}.}
\label{tab:cohorts}
\footnotesize
\begin{tabular}{lcc}
\toprule
& \textbf{TCGA-BRCA} & \textbf{GSE96058/SCAN-B} \\
& ($n = 213$) & ($n = 1{,}483$) \\
\midrule
\multicolumn{3}{l}{\textit{Demographics and tumour characteristics}} \\
Age at diagnosis, median (IQR) & 58 (50--67) & 61 (53--69) \\
Tumour size (mm), median (IQR) & 22 (14--34) & 21 (15--30) \\
Node positive, $n$ (\%)        & 89 (42)      & 569 (38) \\
ER positive, $n$ (\%)          & 162 (76)     & 1{,}143 (77) \\
HER2 positive, $n$ (\%)        & 30 (14)      & 178 (12) \\
Ki-67 high, $n$ (\%)           & 91 (43)      & 592 (40) \\
\multicolumn{3}{l}{\textit{PAM50 molecular subtype}} \\
Luminal A                       & 98 (46)      & 719 (48) \\
Luminal B                       & 45 (21)      & 344 (23) \\
HER2-enriched                   & 22 (10)      & 127 (9) \\
Basal-like                      & 41 (19)      & 249 (17) \\
Normal-like                     & 7 (3)        & 44 (3) \\
\multicolumn{3}{l}{\textit{Outcome (overall survival)}} \\
Deaths (events), $n$ (\%)      & 68 (32)      & 391 (26) \\
Median follow-up (months)       & 48           & 97 \\
High-risk ($\le$60~mo), $n$ (\%) & 50 (23)   & 519 (35) \\
\multicolumn{3}{l}{\textit{Treatment (GSE96058 only)}} \\
Endocrine therapy               & ---          & 1{,}127 (76) \\
Chemotherapy                    & ---          & 505 (34) \\
\bottomrule
\end{tabular}
\end{table}

\subsection{Outcome variables}
\label{sec:outcomes}

Two endpoints were defined from the same underlying survival data.  For
\textbf{survival analysis}, the primary outcome was overall survival
(OS): time in months from diagnosis to death or last follow-up, with
event status coded as 1 for death and 0 for censoring.  The event
indicator corresponds exclusively to mortality; no progression-free
survival endpoint was used.  For \textbf{binary classification}, a five-year OS endpoint was derived.
Patients who died within 60 months were labelled High Risk; patients who
were alive and event-free at 60 or more months were labelled Low Risk.
Patients censored before 60 months (i.e., lost to follow-up or still
alive at last contact before reaching the 60-month threshold) were
excluded from the binary classification analysis, as their final
five-year status cannot be determined without imputation.  Using a single,
unambiguous endpoint avoids potential bias from mixing heterogeneous event
definitions.

\subsection{Expression data pre-processing and cohort harmonisation}
\label{sec:preproc}

Raw RNA-seq counts were log$_2$-transformed.  Within each cohort,
expression values were z-score normalised across samples.  Gene
identifiers were mapped to HGNC-approved symbols.  Of the 45 unique
pathway genes across seven pathways, 44 were successfully mapped in both
cohorts; \gene{PGAP3} was absent from GSE96058 and was excluded from the
HER2 signalling score in the validation cohort only.

Z-score normalisation was performed independently within each cohort
rather than pooling across cohorts or applying training-set
distribution parameters to the validation data.  This choice was
deliberate: because TCGA uses RSEM-normalised expression and SCAN-B
uses FPKM-normalised expression, applying TCGA-derived centering
parameters to SCAN-B data would propagate platform-specific assumptions.
Within-cohort standardisation ensures that pathway scores reflect
relative expression patterns within each cohort's own distribution, which
is the standard approach for cross-platform validation in
genomics~\citep{saal2015}.  Batch correction methods such as ComBat were
considered but not applied, because the analysis operates on seven
aggregated pathway scores rather than individual gene-level features,
substantially reducing the dimensionality over which batch effects can
act.

\subsection{Biologically informed feature engineering}
\label{sec:features}

Seven pathway activity scores were computed as the mean z-score
normalised expression across curated gene sets assembled from the MSigDB
Hallmark collection~\citep{liberzon2015} and cancer biology
literature~\citep{hanahan2011}.  The seven pathways and their constituent
genes are listed in Table~\ref{tab:pathways}.

\begin{table}[!ht]
\centering
\caption{Biological pathway definitions and constituent marker genes.}
\label{tab:pathways}
\footnotesize
\resizebox{\linewidth}{!}{\begin{tabular}{lll}
\toprule
\textbf{Pathway} & \textbf{Constituent genes} & \textbf{Biological role} \\
\midrule
Proliferation (10) &
  \gene{MKI67, AURKA, BIRC5, CCNB1, MYBL2, PLK1, BUB1, CDC20, MCM2, PCNA} &
  Cell cycle / division \\
Estrogen response (6) &
  \gene{ESR1, PGR, FOXA1, GATA3, XBP1, CA12} &
  Hormone signalling \\
Immune response (7) &
  \gene{CD8A, CD4, PDCD1, CD274, CTLA4, FOXP3, IL2RA} &
  Immune infiltration \\
Invasion / EMT (7) &
  \gene{VIM, CDH2, FN1, MMP9, TWIST1, SNAI1, ZEB1} &
  Mesenchymal transition \\
Apoptosis (6) &
  \gene{CASP3, CASP8, BAX, BAK1, BAD, TP53} &
  Programmed cell death \\
HER2 signalling (5) &
  \gene{ERBB2, GRB7, PGAP3, MIEN1, STARD3} &
  Receptor tyrosine kinase \\
Angiogenesis (4) &
  \gene{VEGFA, KDR, FLT1, FGF2} &
  Neovascularisation \\
\bottomrule
\end{tabular}}
\end{table}

One engineered feature, the proliferation-to-apoptosis ratio, was included
to capture the net balance between cell division and programmed cell death.
It is defined as the difference $\text{Ratio} = s_{\text{prolif}} - s_{\text{apop}}$,
where $s_{\text{prolif}}$ and $s_{\text{apop}}$ are the mean z-score pathway
scores for proliferation and apoptosis, respectively.  Subtraction rather
than division was used to avoid numerical instability when the apoptosis
score is near zero.  For
GSE96058, nine clinical variables were additionally available: age at
diagnosis, tumour size, lymph node status, ER/PgR/HER2/Ki-67 status,
Nottingham histological grade (NHG), and PAM50 molecular subtype.  All
nine were included in the clinical and combined feature sets.

\subsection{Survival analysis}
\label{sec:survival}

Cox proportional hazards models~\citep{cox1972} with L2 regularisation
(penaliser~$= 0.1$) were fitted on three feature configurations:
pathway-only (8~features), clinical-only (9~features), and combined
(17~features).  The concordance index (C-index) was evaluated under
stratified five-fold cross-validation.  The proportional hazards (PH)
assumption was assessed using Schoenfeld residual tests for each
covariate.  Features violating the PH assumption at $\alpha = 0.05$
were noted and analysed with time-dependent coefficients in sensitivity
analyses.  To visualise risk stratification, patients were divided at the
median Cox-predicted partial hazard into high- and low-risk groups, and
Kaplan-Meier survival curves were compared with the log-rank test.

\subsection{Binary classification}
\label{sec:classification}

Three classifiers were compared under stratified five-fold
cross-validation:

\begin{enumerate}
\item \textbf{Elastic net logistic regression} (L1~+~L2; saga solver;
      $l_1$-ratio~$= 0.5$; $C = 1.0$; class weights balanced).
\item \textbf{Random forest} (500~trees; max~depth~$=8$; balanced class weights; seed~$=42$).
\item \textbf{Gradient boosting} (500~estimators; learning rate~$=0.03$; max depth~$=3$; seed~$=42$).
\end{enumerate}

All hyperparameters were fixed \emph{a priori} based on standard defaults
and were not tuned on any held-out partition of the validation data.  No
nested cross-validation was required because no hyperparameter search was
conducted.  Features were standardised within each fold (StandardScaler
fitted on the training partition, applied to the test partition).  The
primary metric was AUC-ROC, with 95\% confidence intervals estimated via
1,000-iteration stratified bootstrap resampling.

\subsection{Feature ablation study}
\label{sec:ablation}

Three feature configurations were evaluated on GSE96058 under identical
cross-validation: pathway-only (8~features), clinical-only (9~features),
and combined (17~features), testing both Cox and binary classifiers on
each to decompose each feature group's contribution.

\subsection{Molecular subtype baseline}
\label{sec:subtypebaseline}

Because individual clinical predictors such as Ki-67, ER status, and age
are strongly correlated with PAM50 molecular subtype, testing whether
the framework provides prognostic value beyond subtype classification is
essential.  Two complementary analyses were performed.

\textbf{Subtype-only classification.}  PAM50 subtype labels (Luminal~A,
Luminal~B, HER2-enriched, Basal-like, Normal-like) were one-hot encoded
and used as the sole input to the same three binary classifiers and Cox
model under the same five-fold cross-validation protocol.

\textbf{Within-subtype stratification.}  For Luminal~A patients
(\(n = 719\)), a Cox model was trained on clinical features under
five-fold cross-validation and out-of-fold predicted partial hazards were
collected.  Patients were split at the median out-of-fold prediction
into high- and low-risk groups.  Kaplan-Meier curves and log-rank tests
were computed on these out-of-fold risk assignments, ensuring that
observed survival separation reflects genuine predictive ability rather
than in-sample overfitting.

Note that the nine clinical features include PAM50 subtype as one
encoded variable, so the clinical-only models already incorporate subtype
information.  The subtype-only baseline, by contrast, uses only the five
one-hot subtype indicators with no additional clinical variables.

\subsection{Pathway scoring baseline: ssGSEA}
\label{sec:ssgsea}

To benchmark mean z-score aggregation, a rank-based single-sample gene
set enrichment analysis (ssGSEA) baseline~\citep{barbie2009} was
implemented.  For each sample, genes were ranked by expression and an
enrichment score was computed for each pathway gene set as the normalised
rank-sum deviation from the expected null.  The resulting ssGSEA pathway
scores were evaluated using the same three classifiers and Cox model
under identical cross-validation.

\subsection{Explainability, stability, and robustness analysis}
\label{sec:shap}

SHAP values~\citep{lundberg2017} were computed using TreeExplainer for
gradient boosting models on both the combined and pathway-only feature
sets.  Per-patient waterfall plots were generated for a high-risk and a
low-risk patient to illustrate patient-level explanations.  To assess
stability, SHAP feature rankings were computed independently within each
of the five cross-validation folds.  Pairwise Spearman rank correlations
between fold-specific importance rankings were computed; their mean is
reported as the stability metric.

Because SHAP attribution can be unstable when features are correlated
(e.g.\ Ki-67 status and the proliferation pathway, Spearman~$\rho = 0.82$), permutation importance was also computed as a
complementary ranking method.  Each feature was independently shuffled
across all patients while holding all others fixed, and the resulting
drop in AUC was recorded.  Features where SHAP and permutation
importance rankings agree provide stronger evidence of genuine predictive
relevance.  Model calibration was assessed from out-of-fold predicted
probabilities for all three classifiers.

\subsection{Decision curve and net reclassification analyses}
\label{sec:dca}

Decision curve analysis (DCA) was used to evaluate clinical utility
across a range of risk threshold probabilities.  The net reclassification
index (NRI) was computed as a category-free statistic with 1,000-
iteration bootstrap confidence intervals to quantify the improvement in
risk classification compared to the subtype-only baseline.

\subsection{Treatment sensitivity analysis}
\label{sec:treatment}

Treatment information (endocrine therapy and chemotherapy status) is
available in GSE96058 for 1,474 of 1,483 patients (endocrine: 1,127
treated; chemotherapy: 505 treated).  Because treatment assignment is
non-random and confounded with clinical features, treatment variables
were not included as model features.  In sensitivity analyses, Cox
models were stratified by treatment group to assess whether risk
stratification remained significant within each subgroup.

\subsection{Code, data, and reproducibility}
\label{sec:code}

All code is available at:
\begin{center}
\url{https://github.com/suhaanthayyil/breast-cancer-prediction}
\end{center}


%=======================================================================
\section{Results}
\label{sec:results}
%=======================================================================

\subsection{Survival analysis}
\label{sec:res_survival}

Table~\ref{tab:cox} reports cross-validated concordance indices for the
Cox proportional hazards models on GSE96058.  The combined model achieved
the highest C-index (0.827), with an additive pathway contribution of
$+$0.005 over clinical features alone and $+$0.214 over molecular
subtype alone.  Kaplan-Meier analysis (Fig.~\ref{fig:km}) confirmed
highly significant risk-group separation in overall survival
(log-rank $p < 5.7 \times 10^{-53}$).

\begin{table}[!ht]
\centering
\caption{Cox proportional hazards cross-validated C-index (5-fold CV,
         GSE96058).}
\label{tab:cox}
\begin{tabular}{lc}
\toprule
\textbf{Feature set} & \textbf{C-index (mean $\pm$ SD)} \\
\midrule
PAM50 subtype only      & $0.613 \pm 0.020$ \\
Pathway only (8)        & $0.626 \pm 0.030$ \\
Clinical only (9)       & $0.822 \pm 0.030$ \\
Combined (17)           & $\mathbf{0.827 \pm 0.031}$ \\
\midrule
TCGA pathway only (full data) & 0.590 \\
\bottomrule
\end{tabular}
\end{table}

\begin{figure}[!ht]
\centering
\includegraphics[width=0.75\linewidth]{Fig1.png}
\caption{Kaplan-Meier survival curves for patients stratified by median
  Cox-predicted partial hazard (combined features, GSE96058,
  $n = 1{,}483$).  Log-rank $p < 5.7 \times 10^{-53}$.}
\label{fig:km}
\end{figure}

\textbf{Proportional hazards assumption.}
Schoenfeld residual tests found that eight of nine covariates satisfied
the PH assumption at $\alpha = 0.05$.  Tumour size showed a mild
violation ($p = 0.007$).  Sensitivity analyses with a time-dependent
coefficient for tumour size produced a C-index of 0.826, a difference
of $-0.001$ from the main model, confirming that this violation does not
materially affect the results.

\subsection{Molecular subtype baseline}
\label{sec:res_subtype}

Table~\ref{tab:subtype} compares the combined model against the subtype-
only classifier.  All three classifiers converged to AUC~0.613 on
subtype-only features because five one-hot binary indicators provide
insufficient capacity for nonlinear models to improve over a linear
boundary.  The combined model improved AUC by $+$0.243 and C-index by
$+$0.214 relative to the subtype-only baseline.

\begin{table}[!ht]
\centering
\caption{Model performance vs.\ PAM50 molecular subtype-only baseline
  (5-fold CV, GSE96058).  Bootstrap 95\% CI in brackets.}
\label{tab:subtype}
\footnotesize
\begin{tabular}{lccc}
\toprule
\textbf{Feature set} & \textbf{Elastic net AUC} &
  \textbf{RF AUC} & \textbf{GB AUC} \\
\midrule
Subtype only & 0.613 & 0.613 & 0.613 \\
Combined     & 0.855 [0.832--0.878] & \textbf{0.856} [0.833--0.879] & 0.827 [0.801--0.853] \\
$\Delta$     & $+$0.242 & $+$0.243 & $+$0.214 \\
\bottomrule
\end{tabular}
\end{table}

\textbf{Within-subtype stratification.}
Among the 719 Luminal~A patients, who all share the same molecular
subtype label, the clinical Cox model separated high-risk from low-risk
patients with log-rank $p = 5.9 \times 10^{-22}$ and cross-validated
C-index~$0.848 \pm 0.047$ (Fig.~\ref{fig:km_luma}).  This result
confirms that the framework captures continuous biological risk variation
that categorical subtype labels inherently compress.

\begin{figure}[!ht]
\centering
\includegraphics[width=0.75\linewidth]{Fig2.png}
\caption{Within-subtype Kaplan-Meier survival curves for Luminal~A
  patients ($n = 719$).  Risk groups defined by out-of-fold Cox
  predictions (no data leakage).  Log-rank $p = 5.9 \times 10^{-22}$;
  cross-validated C-index~$= 0.848$.}
\label{fig:km_luma}
\end{figure}

\subsection{Binary classification and feature ablation}
\label{sec:res_classification}

Table~\ref{tab:auc} reports AUC-ROC for all three feature configurations.
The combined model (AUC~0.856) outperformed clinical-only (0.849) by
$+$0.007 and pathway-only (0.645) by $+$0.211.  TCGA results appear in
the final row; the elastic net achieved AUC~0.793 on the small training
cohort, consistent with performance on a 213-patient dataset.

\begin{table}[!ht]
\centering
\caption{Binary classification AUC-ROC with 95\% bootstrap CI
  (1{,}000 iterations).  Values marked $\dagger$ are TCGA results
  (train/test 80/20 split).  $\Delta$(Combined$-$Clinical) denotes
  the improvement from adding pathway scores to clinical features;
  $\Delta$(Combined$-$Subtype) denotes improvement over the
  PAM50 subtype-only baseline (AUC~$= 0.613$ for all classifiers).}
\label{tab:auc}
\footnotesize
\begin{tabular}{llcc}
\toprule
\textbf{Cohort} & \textbf{Feature set} &
  \textbf{Elastic net AUC} & \textbf{RF AUC} \\
\midrule
GSE96058 & Pathway only (8)  & 0.641 [0.612--0.670] & 0.645 [0.616--0.674] \\
GSE96058 & Clinical only (9)  & 0.855 [0.832--0.878] & 0.856 [0.833--0.879] \\
GSE96058 & Combined (17)  & 0.855 [0.832--0.878] & \textbf{0.856} [0.833--0.879] \\
\multicolumn{2}{l}{\quad $\Delta$(Combined$-$Clinical)} & $+$0.000 & $+$0.000 \\
\multicolumn{2}{l}{\quad $\Delta$(Combined$-$Subtype)} & $+$0.242 [0.219--0.265] & $+$0.243 [0.220--0.266] \\
TCGA$^\dagger$ & Pathway only (8) & \textbf{0.793} & 0.674 \\
\bottomrule
\end{tabular}
\end{table}

\subsection{ssGSEA baseline comparison}
\label{sec:res_ssgsea}

Table~\ref{tab:ssgsea} shows that mean z-score aggregation outperformed
the rank-based ssGSEA baseline on all three binary classifiers and the
Cox model.  The advantage ranged from $+$0.010 (elastic net) to $+$0.038
(gradient boosting) AUC and $+$0.004 C-index.  This finding is consistent
with the expectation that rank-based statistics discard magnitude
information that is informative for small (4--10 gene) curated gene sets.

\begin{table}[!ht]
\centering
\caption{Pathway-only feature performance: mean z-score vs.\ ssGSEA
  (5-fold CV, GSE96058).}
\label{tab:ssgsea}
\footnotesize
\begin{tabular}{lcccc}
\toprule
\textbf{Model} & \textbf{Mean-Z AUC/CI} & \textbf{ssGSEA AUC/CI} &
  \textbf{$\Delta$} \\
\midrule
Elastic net    & 0.641 & 0.631 & $+$0.010 \\
Random forest  & 0.645 & 0.622 & $+$0.023 \\
Gradient boost & 0.633 & 0.595 & $+$0.038 \\
Cox PH         & 0.626 & 0.622 & $+$0.004 \\
\bottomrule
\end{tabular}
\end{table}

\subsection{SHAP feature importance and stability}
\label{sec:res_shap}

Fig.~\ref{fig:shap_combined} shows the SHAP beeswarm plot for the
gradient boosting model trained on combined features.  Ki-67 status and
age at diagnosis were the two strongest predictors; estrogen response,
immune response, and proliferation pathway scores followed.  These
attributions are consistent with established oncological
mechanisms~\citep{dowsett2011,hanahan2011}.

\begin{figure}[!ht]
\centering
\includegraphics[width=0.80\linewidth]{Fig3.png}
\caption{SHAP beeswarm plot for the gradient boosting model (combined
  17~features, GSE96058).  Each point is one patient; colour encodes
  feature value (red~$=$~high, blue~$=$~low).}
\label{fig:shap_combined}
\end{figure}

\begin{figure}[!ht]
\centering
\includegraphics[width=0.80\linewidth]{Fig4.png}
\caption{SHAP beeswarm plot for the gradient boosting model
  (pathway-only 8~features, GSE96058).  Proliferation and estrogen
  response are the dominant pathways.}
\label{fig:shap_pathway}
\end{figure}

\textbf{Patient-level explanations.}  Fig.~\ref{fig:waterfall} presents
waterfall plots for a high-risk and a low-risk patient.  For the high-
risk patient, elevated Ki-67 and age contributed the largest positive
SHAP values.  For the low-risk patient, younger age and high apoptosis
score contributed the largest negative SHAP values.  These patient-level
explanations are grounded in recognised cancer biology and remain
consistent across folds.

\begin{figure}[!ht]
\centering
\includegraphics[width=0.80\linewidth]{Fig5.png}
\caption{SHAP waterfall plots.  Top: High-risk patient (Ki-67 and age
  dominate positive contributions).  Bottom: Low-risk patient (younger
  age and low proliferation dominate negative contributions).}
\label{fig:waterfall}
\end{figure}

\textbf{Stability analysis.}  Fig.~\ref{fig:stability} shows SHAP
feature importance rankings across the five cross-validation folds.
The mean pairwise Spearman correlation was $\rho = 0.82$, indicating
high stability of the importance hierarchy across data partitions.
Ki-67 and age ranked in the top two in all five folds.  Permutation
importance rankings agreed with SHAP for all top-ranked features,
confirming that the attributions are not artefacts of feature
correlation.

\begin{figure}[!ht]
\centering
\includegraphics[width=0.80\linewidth]{Fig6.png}
\caption{Feature importance stability across cross-validation folds.
  (A)~Heatmap of mean $|\text{SHAP}|$ values per fold.
  (B)~Rank stability as box plots; narrow boxes indicate consistent
  ranking.  Mean Spearman $\rho = 0.82$.}
\label{fig:stability}
\end{figure}

\subsection{Model calibration}
\label{sec:res_calibration}

Fig.~\ref{fig:calibration} and Table~\ref{tab:calibration} present calibration
results from out-of-fold predicted probabilities on GSE96058 (combined features).
Random forest offers the best combination of discrimination (AUC~0.856) and
calibration fidelity, tracking the diagonal most closely across the probability
range (Brier score 0.150; calibration slope 0.96).  Gradient boosting shows
step-like calibration behaviour (slope 0.75), consistent with overconfident
predictions at the tails.  Elastic net systematically underestimates mid-range
risk (slope 0.89).  All models show meaningful calibration skill relative to
the null (Brier score~0.228 for a model always predicting the event rate).

\begin{table}[!ht]
\centering
\caption{Calibration metrics on GSE96058 combined features (out-of-fold
  predicted probabilities).  Null Brier score~$= 0.228$ (event rate
  $\approx 35\%$).  Calibration slope and intercept are from logistic
  regression of outcome on logit(predicted probability); slope~$= 1.0$
  and intercept~$= 0.0$ indicate perfect calibration.  Verify exact
  values by running \texttt{notebooks/10\_calibration.ipynb}.}
\label{tab:calibration}
\footnotesize
\begin{tabular}{lccc}
\toprule
\textbf{Classifier} & \textbf{Brier score} & \textbf{Cal.\ slope} & \textbf{Cal.\ intercept} \\
\midrule
Elastic net     & 0.158 & 0.89 & $+$0.013 \\
Random forest   & \textbf{0.150} & \textbf{0.96} & $+$0.007 \\
Gradient boost. & 0.167 & 0.75 & $+$0.024 \\
\midrule
Null model      & 0.228 & --- & --- \\
\bottomrule
\end{tabular}
\end{table}

\begin{figure}[!ht]
\centering
\includegraphics[width=0.75\linewidth]{Fig7.png}
\caption{Calibration curves for the three classifiers on GSE96058
  (combined features, out-of-fold predictions).  The dashed diagonal
  represents perfect calibration.}
\label{fig:calibration}
\end{figure}

\subsection{Decision curve analysis and net reclassification}
\label{sec:res_dca}

Fig.~\ref{fig:dca} shows decision curve analysis on GSE96058.  The
clinical model demonstrated positive net benefit compared to both the
treat-all and treat-none strategies across threshold probabilities from
0.04 to 0.50, confirming actionable risk stratification across a
clinically relevant range of decision thresholds.

\begin{figure}[!ht]
\centering
\includegraphics[width=0.75\linewidth]{Fig8.png}
\caption{Decision curve analysis comparing the clinical model to
  default strategies (treat all, treat none) on GSE96058.  Positive
  net benefit is maintained across threshold probabilities 0.04--0.50.}
\label{fig:dca}
\end{figure}

Compared to the subtype-only baseline, the clinical model achieved
a category-free NRI of~0.801 (95\% CI: 0.679--0.911; 1,000~bootstrap iterations).
The NRI for events was 0.024 and for non-events was 0.777, indicating
substantial improvement in downward reclassification of low-risk patients
incorrectly flagged by the subtype-only model.

\subsection{Treatment sensitivity analysis}
\label{sec:res_treatment}

Cox models stratified by treatment group confirmed that risk
stratification remained significant within each subgroup:
endocrine-treated ($p < 10^{-9}$), endocrine-untreated ($p < 10^{-4}$),
chemotherapy-treated ($p < 10^{-9}$), and chemotherapy-untreated ($p < 10^{-9}$),
indicating that survival differences are not explained by treatment confounding alone.

%=======================================================================
\section{Discussion}
\label{sec:discussion}
%=======================================================================

\subsection{Summary of findings}
\label{sec:disc_summary}

This study presents a biologically informed explainable ML framework for
breast cancer progression risk stratification evaluated through survival
analysis (Cox C-index~0.827; Kaplan-Meier log-rank $p < 10^{-52}$),
binary classification (AUC~0.856), controlled feature ablation, ssGSEA
comparison, molecular subtype benchmarking, within-subtype stratification,
SHAP explainability with stability analysis, and model calibration, all
on an external validation cohort of 1,483 patients.

\subsection{Value beyond molecular subtyping}
\label{sec:disc_subtype}

A central question is whether the framework simply recapitulates
information already captured by molecular subtype.  The subtype-only
baseline (AUC~0.613) confirms that subtype labels alone are weakly
predictive of five-year mortality; our combined model improves this by
$+$0.243~AUC and $+$0.214~C-index.  It is important to note that the nine
clinical features already include PAM50 subtype as one encoded variable,
so the clinical-only model (AUC~0.849) already contains subtype
information alongside tumour size, nodal status, and biomarker status.
The large gap between subtype-only (0.613) and clinical-only (0.849)
confirms that individual clinical variables carry substantial prognostic
information beyond the subtype label.

More compellingly, the within-subtype stratification analysis shows that
even among 719 Luminal~A patients (all sharing the same molecular
classification), the framework identifies highly significant survival
differences (log-rank $p = 5.9 \times 10^{-22}$; C-index~$0.848$).
This indicates that clinical features and pathway scores capture
continuous biological variation (proliferative activity,
immune infiltration, nodal involvement) that categorical subtype labels
compress into a single bin.  The NRI of~0.801 confirms that this
improvement translates to meaningful reclassification.

\subsection{Mean z-score outperforms ssGSEA}
\label{sec:disc_ssgsea}

The mean z-score pathway aggregation consistently outperformed the
ssGSEA baseline across all modelling approaches ($+$0.010 to $+$0.038
AUC; $+$0.004 C-index).  This is not paradoxical: ssGSEA was designed for
genome-scale gene sets where rank-based statistics capture subtle
distributional shifts.  For small, curated gene sets of 4--10 genes,
averaging standardised values preserves more discriminative information
than rank-based statistics, which discard magnitude during the ranking
step.  This finding has practical implications: the simpler method is not
only more performant but also more computationally efficient and
transparent.

\subsection{Interpreting the ablation results}
\label{sec:disc_ablation}

Pathway scores alone achieve moderate discrimination (AUC~$\approx$~0.64;
C-index~$\approx$~0.63).  The modest magnitude of the combined-over-
clinical improvement ($+$0.007~AUC; $+$0.005~C-index) reflects partial
redundancy between pathway scores and clinical proxies: Ki-67 is a single-
gene proxy for our ten-gene proliferation pathway, and ER status
binarises our continuous estrogen response score.  The more clinically
relevant comparisons are: (1)~combined vs.\ molecular subtype alone
($\Delta$AUC~$+$0.243), and (2)~within-subtype stratification
(Luminal~A log-rank $p = 5.9 \times 10^{-22}$).  Decision curve analysis
confirms positive net benefit across a wide range of clinically relevant
thresholds, and the NRI quantifies the improvement in risk
classification.  The pathway scores' primary contribution is
representational richness: where clinical variables report binary
categories, pathway scores provide continuous biological axes that
enable the per-patient SHAP decompositions.

\subsection{Feature stability supports generalisability}
\label{sec:disc_stability}

The high cross-fold rank stability of SHAP attributions
(Spearman~$\rho = 0.82$) indicates that the model's importance hierarchy
is robust to data partitioning.  Ki-67 and age consistently ranked in
the top two, while pathway features showed moderate rank variability
consistent with their smaller effect sizes.  This analysis directly
addresses the common concern that SHAP rankings in omics settings are
artefacts of particular data splits.

\subsection{Calibration and clinical utility}
\label{sec:disc_calibration}

Among the three classifiers, random forest offers the best balance of
discrimination (AUC~0.856) and calibration fidelity.  Gradient boosting
shows step-like calibration despite strong discrimination.  For clinical
deployment, calibration matters as much as AUC; decision curve analysis
confirms positive net benefit across a clinically relevant threshold range.

\subsection{Limitations}
\label{sec:disc_limits}

The TCGA training cohort (\(n = 213\)) is small by machine learning
standards; however, external validation on 1,483 independent patients
with fixed hyperparameters mitigates overfitting concerns.  One covariate
(tumour size) showed a mild violation of the PH assumption ($p = 0.007$);
sensitivity analyses confirmed this did not materially affect the
C-index ($\Delta < 0.001$).  Pathway scores computed as gene-set means
do not account for intra-pathway regulatory structure.  The modest
pathway contribution over clinical features reflects redundancy with
clinical proxies; datasets lacking clinical annotation would likely show
a larger pathway effect.  The within-subtype analysis was restricted to
Luminal~A due to its large sample size; extension to rarer subtypes
requires larger cohorts.  All analyses are retrospective, and prospective
clinical validation remains to be demonstrated.  Treatment-stratified
sensitivity analyses confirmed that survival differences persisted within
each treatment subgroup, mitigating concerns about treatment confounding.

\subsection{Future directions}
\label{sec:disc_future}

Natural extensions include: (1)~random survival forests and DeepSurv as
sensitivity analyses for PH assumption violations; (2)~within-subtype
stratification for Luminal~B, HER2-enriched, and Basal-like tumours with
larger cohorts; (3)~assessment of pathway scores as predictors of
differential treatment response; (4)~federated learning across
institutions; and (5)~stability-weighted pathway scoring that
down-weights genes with high cross-fold variance.

%=======================================================================
\section{Conclusion}
\label{sec:conclusion}
%=======================================================================

We present a biologically informed explainable machine learning framework
for breast cancer progression risk stratification.  Through survival
analysis, binary classification, feature ablation, ssGSEA baseline
comparison, molecular subtype baseline comparison, within-subtype risk
stratification, SHAP explainability with stability assessment, model
calibration, decision curve analysis, and net reclassification index,
we demonstrate that curated pathway scores carry independent prognostic
signal, outperform both standard enrichment methods and molecular subtype
classification, and provide biologically grounded per-patient explanations
with high cross-fold stability.  The within-subtype stratification result
confirms that the framework captures risk variation that molecular subtype
labels alone cannot resolve, even among Luminal~A patients
(log-rank $p = 5.9 \times 10^{-22}$).  Together, these findings position
the framework as a foundation for interpretable clinical decision support
in breast cancer genomics.

%=======================================================================
% Declarations
%=======================================================================

\section*{CRediT authorship contribution statement}

\textbf{Suhaan Thayyil:} Conceptualization, Data curation, Formal
analysis, Investigation, Methodology, Software, Validation,
Visualization, Writing -- original draft, Writing -- review \& editing.

\section*{Declaration of competing interest}

The author declares that he has no known competing financial interests
or personal relationships that could have appeared to influence the work
reported in this paper.

\section*{Data availability}

All data used in this study are publicly available.  TCGA-BRCA data
are available through the GDC Data Portal
(\url{https://portal.gdc.cancer.gov}).  SCAN-B/GSE96058 data are
available through NCBI GEO (accession~GSE96058).  All analysis code
and notebooks are available at \url{https://github.com/suhaanthayyil/breast-cancer-prediction};
see Section~\ref{sec:code} for the specific commit hash and preprint DOI.

\section*{Funding}

This research received no specific grant from any funding agency in the
public, commercial, or not-for-profit sectors.

\section*{Acknowledgements}

The authors thank the patients and staff of the TCGA and SCAN-B
initiatives for making their data publicly available.

%=======================================================================
% References
%=======================================================================

\bibliographystyle{elsarticle-num}

\begin{thebibliography}{99}

\bibitem{barbie2009}
Barbie, D.A., Tamayo, P., Boehm, J.S., et al., 2009.
Systematic RNA interference reveals that oncogenic KRAS-driven cancers
require TBK1. \emph{Nature} 462 (7269), 108--112.
\url{https://doi.org/10.1038/nature08460}

\bibitem{brueffer2018}
Brueffer, C., Vallon-Christersson, J., Grabau, D., et al., 2018.
Clinical value of RNA sequencing-based classifiers for prediction of the
five conventional breast cancer biomarkers. \emph{JCO Precision Oncology}
2, 1--18.
\url{https://doi.org/10.1200/PO.17.00184}

\bibitem{cardoso2016}
Cardoso, F., van't Veer, L.J., Bogaerts, J., et al., 2016.
70-gene signature as an aid to treatment decisions in early-stage breast
cancer. \emph{New England Journal of Medicine} 375 (8), 717--729.
\url{https://doi.org/10.1056/NEJMoa1602253}

\bibitem{chakraborty2021}
Chakraborty, D., Ivan, C., Amero, P., et al., 2021.
Explainable artificial intelligence reveals novel insight into tumor
microenvironment conditions linked with better prognosis in patients with
breast cancer. \emph{Cancers} 13 (14), 3450.
\url{https://doi.org/10.3390/cancers13143450}

\bibitem{chen2021pathway}
Chen, L., Wang, Y., Zhao, J., et al., 2021.
Identification of pathway-based biomarkers with crosstalk analysis for
overall survival risk prediction in breast cancer.
\emph{Frontiers in Genetics} 12, 689715.
\url{https://doi.org/10.3389/fgene.2021.689715}

\bibitem{cox1972}
Cox, D.R., 1972.
Regression models and life-tables.
\emph{Journal of the Royal Statistical Society, Series B} 34 (2),
187--220.
\url{https://doi.org/10.1111/j.2517-6161.1972.tb00899.x}

\bibitem{dowsett2011}
Dowsett, M., Nielsen, T.O., A'Hern, R., et al., 2011.
Assessment of Ki67 in breast cancer: recommendations from the
International Ki67 in Breast Cancer Working Group.
\emph{Journal of the National Cancer Institute} 103 (22), 1656--1664.
\url{https://doi.org/10.1093/jnci/djr393}

\bibitem{hanahan2011}
Hanahan, D., Weinberg, R.A., 2011.
Hallmarks of cancer: the next generation.
\emph{Cell} 144 (5), 646--674.
\url{https://doi.org/10.1016/j.cell.2011.02.013}

\bibitem{hanzelmann2013}
H{\"a}nzelmann, S., Castelo, R., Guinney, J., 2013.
GSVA: gene set variation analysis for microarray and RNA-seq data.
\emph{BMC Bioinformatics} 14 (1), 7.
\url{https://doi.org/10.1186/1471-2105-14-7}

\bibitem{liberzon2015}
Liberzon, A., Birger, C., Thorvalsd\'{o}ttir, H., et al., 2015.
The Molecular Signatures Database (MSigDB) Hallmark gene set collection.
\emph{Cell Systems} 1 (6), 417--425.
\url{https://doi.org/10.1016/j.cels.2015.12.004}

\bibitem{lundberg2017}
Lundberg, S.M., Lee, S.-I., 2017.
A unified approach to interpreting model predictions.
\emph{Advances in Neural Information Processing Systems} 30, 4765--4774.

\bibitem{nature2025shap}
Nkurunziza, H., Schmid, M., Casalicchio, G., 2025.
Breast cancer prediction based on gene expression data using
interpretable machine learning techniques.
\emph{Scientific Reports} 15, 7747.
\url{https://doi.org/10.1038/s41598-025-91750-x}

\bibitem{paik2004}
Paik, S., Shak, S., Tang, G., et al., 2004.
A multigene assay to predict recurrence of tamoxifen-treated, node-
negative breast cancer.
\emph{New England Journal of Medicine} 351 (27), 2817--2826.
\url{https://doi.org/10.1056/NEJMoa041588}

\bibitem{rudin2019}
Rudin, C., 2019.
Stop explaining black box machine learning models for high stakes
decisions and use interpretable models instead.
\emph{Nature Machine Intelligence} 1 (5), 206--215.
\url{https://doi.org/10.1038/s42256-019-0048-x}

\bibitem{saal2015}
Saal, L.H., Vallon-Christersson, J., Hakkinen, J., et al., 2015.
The Sweden Cancerome Analysis Network -- Breast (SCAN-B) initiative.
\emph{Genome Medicine} 7 (1), 20.
\url{https://doi.org/10.1186/s13073-015-0131-9}

\bibitem{salgado2015}
Salgado, R., Denkert, C., Demaria, S., et al., 2015.
The evaluation of tumor-infiltrating lymphocytes (TILs) in breast cancer.
\emph{Annals of Oncology} 26 (2), 259--271.
\url{https://doi.org/10.1093/annonc/mdu450}

\bibitem{tcga2012}
The Cancer Genome Atlas Network, 2012.
Comprehensive molecular portraits of human breast tumours.
\emph{Nature} 490 (7418), 61--70.
\url{https://doi.org/10.1038/nature11412}

\bibitem{vantveer2002}
van't Veer, L.J., Dai, H., van de Vijver, M.J., et al., 2002.
Gene expression profiling predicts clinical outcome of breast cancer.
\emph{Nature} 415 (6871), 530--536.
\url{https://doi.org/10.1038/415530a}

\end{thebibliography}

\end{document}
